% GENERATED from the paper's sources; edit those, not this file.
\begin{description}
  \item[Clock] Read as clock rates, the weights say how fast time runs at each place, set by the saddle rather than chosen place by place, and a place that falls short is a place whose clock has stopped.~\cite{PageWootters1983} \emph{(reading)} It reads the rate of a test clock's phase at each place per unit of the clock's reading, in units of its frequency, under premises~\ref{premise:evolution}, \ref{premise:probes}. Every simulated clock runs at its place's weight, to within $10^{-9}$. On the ring of $6$ qubits the first cell's clock runs at $0.015$ at depth $0.2$, and past its drop-out, at depth $1$, it stops; on the even XXZ ring of $4$ cells the deepened cell's clock still runs at depth $8$, at $0.0081$.
  \item[Chance from records] What a part reads as chance is which record its description is conditioned on: in the whole's frame nothing is drawn, and the readings can be correlated beyond what one description of absolute outcomes allows only while no copy of them outside their places is full.~\cite{BongEtAl2020} \emph{(reading)} It reads the CHSH value of Alice's and Bob's outcomes against the overlap of the outside copy of the friends' records, under premises~\ref{premise:evolution}, \ref{premise:records}. Solved once with no time parameter, the protocol's CHSH value is $(1+\gamma)^2/\sqrt{2}$ to within $10^{-9}$, from $2.828$ with no copy to $0.7071$ with a full one, above the bound $2$, with these settings, only while the overlap exceeds $0.6818$ (with turned settings, derived rather than simulated, at every overlap short of a full copy); nothing signals, to the bound $10^{-12}$.
  \item[Lapse] The weights hold a static lapse's place: they weight the generator under which the saddle state is stationary, turn each place's energy into the common value at their own rate, as a lapse turns local energy into energy measured from far away, and vanish where a place's energy no longer counts.~\cite{Tolman1930} \emph{(matched)} It reads the energy an excitation at a place adds to the weighted generator, per unit of its own energy, under premises~\ref{premise:evolution}, \ref{premise:probes}. Exciting the probe at each place of the ring adds its place's weight times the probe's energy, to within $10^{-12}$ of the weight.
  \item[Tolman temperature] In equilibrium a place's temperature times its weight is one constant across the system, as temperature times the lapse is in a static gravitational field, so a place with less weight runs hotter; as a common bound on every local energy falls to the common value, the coldest equilibrium it allows has its local inverse temperatures approach the saddle weights' proportions.~\cite{Tolman1930} \emph{(matched)} It reads the temperature a test thermometer reads at each place, in units of the global temperature, under premises~\ref{premise:evolution}, \ref{premise:probes}, \ref{premise:equilibrium}. Every simulated thermometer reads a temperature times weight equal to the global temperature, to within $10^{-9}$; at depth $0.2$ the first cell runs $11$ times hotter than a cell of the uniform ring, and past its drop-out its thermometer has no finite reading.
  \item[Proper length] A cell's proper length is its radar length, the speed of light times the time long waves take to cross it, read by the clock at the cell: every cell has the same proper length whatever its weight, and the weights enter the geometry only through the lapse.~\cite{DubailEtAl2017} \emph{(matched)} It reads the time a long wave takes to cross a stretch, read by the clocks along it, and the ground state's entanglement across a cut, under premises~\ref{premise:evolution}, \ref{premise:probes}, \ref{premise:cone}. Long waves on the weighted chain take the optical time to within the relative bound $0.001$, and the saddle state's own entanglement is the vacuum's on the metric whose lapse is its weights, with a cutoff of one cell, to within $10^{-4}$.
  \item[Horizon] A place with no weight reads as a place where a static lapse vanishes, as it does at a Killing horizon; here it sits wherever a place's density falls below the common value.~\cite{Gourgoulhon2007} \emph{(reading)} It reads at a place with no weight, its clock's rate, its thermometer's reading, and the energy its excitation adds to the whole, under premises~\ref{premise:evolution}, \ref{premise:probes}, \ref{premise:equilibrium}. Past its drop-out the first cell's simulated clock stops, its thermometer has no finite reading, and an excitation there adds nothing to the weighted generator's energy.
  \item[Cosmological constant] The part of the common value set by the whole's size reads as the negative of a cosmological constant: a value the whole sets, not any one place's terms.~\cite{HenneauxTeitelboim1989} \emph{(reading)} It reads the part of the common value set by the size of a uniform critical ring, and the central charge it gives, under premises~\ref{premise:constant}. On uniform critical rings of $8$ to $60$ cells the size-set part gives a central charge $0.5000$, and an interval's entanglement entropy gives $0.5003$ independently.
  \item[Horizon temperature] A part's edges read as its horizons: the weights of its entanglement Hamiltonian vanish there, and the temperature their slope gives, the speed times the slope over 2 pi, equals the local temperature at the part's centre, as Tolman's relation across a static patch requires.~\cite{BisognanoWichmann1976} \emph{(matched)} It reads the ratio of a part's horizon temperature, the speed times its weights' slope at the edge over 2 pi, to the local temperature at its centre, under premises~\ref{premise:subsystem}, \ref{premise:cone}. For intervals of $8$ to $40$ cells of the critical chain, the ratio read with the weights' longer-range terms differs from one by $1.4\times10^{-4}$ to $5.6\times10^{-6}$; from the nearest terms alone it is $1.05$.
  \item[De Sitter temperature] Read as the cosmological horizon, a part's horizon has the surface gravity general relativity gives the cosmological constant, and the temperature the part's weights give it is Gibbons and Hawking's.~\cite{GibbonsHawking1977} \emph{(reading)} It reads the temperature of the cosmological horizon, for the measured cosmological constant, under premises~\ref{premise:subsystem}, \ref{premise:cone}, \ref{premise:patch}. Not simulated: the system supplies the 2 pi. With the measured Hubble rate $67.36$ km/s/Mpc and dark-energy fraction $0.6847$, the temperature is $2.20\times10^{-30}\,\mathrm{K}$, with the $1\%$ uncertainty those values carry.
  \item[Shadow price] Each weight is the shadow price of its place's bound: the rate at which the common value moves when that bound is relaxed.~\cite{Kantorovich1960} \emph{(matched)} It reads the rate at which the common value moves when one place's bound is relaxed. Relaxing and tightening each cell's bound on the ring and solving again moves the common value, between the two, at that cell's weight, to within $10^{-4}$.
\end{description}
